\ea\label{ex:text18-1}
O sinivi potee\\
\gll o	sinivi	potee\\
     \textsc{art}3.\textsc{sg}	canoe	potee.canoe\\
\glt `The \textit{potee} canoe'
\z
\largerpage[2]

\ea\label{ex:text18-2}
O potee o meho kaku sinivi to tei nana Teapu.\\
\gll o	potee	o	meho	kaku	sinivi	to	tei	nana	teapu\\
     \textsc{art}3.\textsc{sg}	potee.canoe	\textsc{art}3.\textsc{sg}	other3.\textsc{sg}	type	canoe	\textsc{rel}	be.there	3\textsc{sg}.\textsc{ipfv}	Teop.area\\
\glt `The \textit{potee} canoe is another type of canoe in the Teop area.'\footnote{lit. `... another type of canoe that is in the Teop region.'}
\z

\ea\label{ex:text18-3}
Enaa kahi vahutate nom, havee to pakupaku vaavaha me raara e.\\
\gll e-naa	kahi	vahutate	nom	havee	to	paku{$\sim$}paku	vaavaha	me	ra=ara	e\\
     \textsc{h}-1\textsc{sg}.\textsc{pron}	\textsc{tm}1	tell	\textsc{ipfv}	what	\textsc{rel}	\textsc{red}{$\sim$}make	different	with	1\textsc{pl}.\textsc{in}.\textsc{ipfv}=1\textsc{pl}.\textsc{in}.\textsc{pron}	3\textsc{sg}.\textsc{pron}\\
\glt `I am going to tell you how we make it.'\footnote{lit. `... what (it is) like that we differently make it with.'}
\z

\ea\label{ex:text18-4}
Beara rake tea paku o potee, eara toro tagihu roho rasuu tea tara rakerake to naono vai to mataa nana tea paku bono potee.\\
\gll be=ara	rake	tea	paku	o	potee	e-ara	toro	tagihu	roho	rasuu	tea	tara	rake{$\sim$}rake	to	naono	vai	to	mataa	nana	tea	paku	bono	potee\\
     when=1\textsc{pl}.\textsc{in}.\textsc{pron}	want	\textsc{compl}1	make	\textsc{art}3.\textsc{sg}	potee	\textsc{h}-1\textsc{pl}.\textsc{in}.\textsc{pron}	\textsc{tm}7	go.into.jungle	before	jungle	\textsc{compl}1	look	\textsc{red}{$\sim$}search	\textsc{nspec}3.\textsc{sg}	tree	\textsc{dem}6	\textsc{nspec}3.\textsc{sg}	good	3\textsc{sg}.\textsc{ipfv}	\textsc{compl}1	make	\textsc{obj}.\textsc{art}3.\textsc{sg}	potee.canoe\\
\glt `When we want to make a \textit{potee} canoe, we must first go into the jungle to look for a tree that is good for making the potee.'
\z

\ea\label{ex:text18-5}
A kiu no potee sa beera potee haana bona kiu vaa teo sinivi suuna,\\
\gll a	kiu	n=o	potee	sa	beera	potee	haa=na	bona	kiu	vaa	te=o	sinivi	suuna\\
     \textsc{art}2.\textsc{sg}	work	3\textsc{sg}.\textsc{poss}=\textsc{art}3.\textsc{sg}	potee.canoe	\textsc{neg}1	big	be.like	\textsc{neg}2=3\textsc{sg}.\textsc{ipfv}	\textsc{obj}.\textsc{art}2.\textsc{sg}	work	\textsc{lnk}	\textsc{prep}=\textsc{art}3.\textsc{sg}	canoe	indigenous\\
\glt `The work of (making) the \textit{potee} canoe is not as big as the work on the indigenous canoe,'
\z\clearpage

\ea\label{ex:text18-6}
ei mea peha aba na antee nana tea paku bona.\\
\gll ei	me=a	peha	aba	na	antee	nana	tea	paku	bona\\
     \textsc{dem}2	and4=\textsc{art}2.\textsc{sg}	one2.\textsc{sg}	person	\textsc{tm}2	can	3\textsc{sg}.\textsc{ipfv}	\textsc{compl}1	do	4\textsc{sg}.\textsc{pron}\\
\glt `that is, a single person can do it.'
\z

\ea\label{ex:text18-7}
Me bea peha aba rova koa vurahe to naono, eove repaa tagava gunaha bona.\\
\gll me	be=a	peha	aba	rova	koa	vurahe	to	naono	eove	re=paa	tagava	gunaha	bona\\
     and4	when=\textsc{art}2.\textsc{sg}	one2.\textsc{sg}	person	even	just	find	\textsc{nspec}3.\textsc{sg}	tree	\textsc{h}.3\textsc{sg}.\textsc{pron}	\textsc{consec}=\textsc{tm}3	fell	go.down	4\textsc{sg}.\textsc{pron}\\
\glt `And even when just a single person finds a (suitable) tree, he fells it.'
\z

\ea\label{ex:text18-8}
Be tagava bona, eove repaa rom kurus bono govuhu no naono, be tara voen na mataa nana.\\
\gll be	tagava	bona	eove	re=paa	rom	kurus	bono	govuhu	n=o	naono	be	tara	vo=en	na	mataa	nana\\
     when	fell	4\textsc{sg}.\textsc{pron}	\textsc{h}.3\textsc{sg}.\textsc{pron}	\textsc{consec}=\textsc{tm}3	cut.crosswise	cut.off	\textsc{obj}.\textsc{art}3.\textsc{sg}	crown	3\textsc{sg}.\textsc{poss}=\textsc{art}3.\textsc{sg}	tree	if	see	like=\textsc{dem}3	\textsc{tm}2	good	3\textsc{sg}.\textsc{ipfv}\\
\glt `When (he) has felled it, he cuts the crown of the tree off, if he sees that (the log) is good.'
\z

\ea\label{ex:text18-9}
Beve rom kurusu vakavara bona, eove repaa poana pete bona.\\
\gll beve	rom	kurusu	va-kavara	bona	eove	re=paa	poana	pete	bona\\
     when.3\textsc{sg}.\textsc{pron}	cut.crosswise	cut.off	\textsc{advr}-complete	4\textsc{sg}.\textsc{pron}	\textsc{h}.3\textsc{sg}.\textsc{pron}	\textsc{consec}=\textsc{tm}3	slice.upper.side.off	subsequently	4\textsc{sg}.\textsc{pron}\\
\glt `When he has cut it (the crown) off, he slices the upper side of it (the log).'
\z

\ea\label{ex:text18-10}
A kiu vaa teo potee a varuruva koa tabae o koma beera.\\
\gll a	kiu	vaa	te=o	potee	o	varuruva	koa	tabae	o	koma	beera\\
     \textsc{art}2.\textsc{sg}	work	\textsc{lnk}	\textsc{prep}=\textsc{art}3.\textsc{sg}	potee.canoe	\textsc{art}3.\textsc{sg}	easy	just	because	\textsc{art}3.\textsc{sg}	inside	big\\
\glt `The work on the \textit{potee} canoe is easy, because (it has) a big inside.'
\z

\ea\label{ex:text18-11}
Be poana vakavara bona, eove repaa koro bona bona toraaraa ae a rom.\\
\gll be	poana	va-kavara	bona	eove	re=paa	koro	bona	bona	toraara	ae	a	rom\\
     when	slice.upper.part.off	\textsc{advr}-complete	4\textsc{sg}.\textsc{pron}	\textsc{h}.3\textsc{sg}.\textsc{pron}	\textsc{consec}=\textsc{tm}3	carve.out	4\textsc{sg}.\textsc{pron}	\textsc{art}2.\textsc{sg}	axe	and1	\textsc{art}2.sg	hoe\\
\glt `When having sliced off its upper part, he carves the inside with an axe and a hoe.'
\z

\ea\label{ex:text18-12}
A kiu vaa tea koro koma sinivi a beera,\\
\gll a	kiu	vaa	tea	koro	koma	sinivi	a	beera\\
     \textsc{art}2.\textsc{sg}	work	\textsc{lnk}	\textsc{compl}1	carve	inside	indigenous.canoe	\textsc{art}2.\textsc{sg}	big\\
\glt `The work of carving the inside of the canoe is big,'
\z

\ea\label{ex:text18-13}
me be nagunagu, eove repaa kaokahi bona ore paa vaavuriha.\\
\gll me	be	nagu{$\sim$}nagu	eove	re=paa	kao-kahi	bona	o=re	paa	vaavuriha\\
     and4	when	\textsc{red}{$\sim$}exhausted	\textsc{h}.3\textsc{sg}.\textsc{pron}	\textsc{consec}=\textsc{tm}3	go-from	4\textsc{sg}.\textsc{pron}	3.\textsc{pron}=\textsc{consec}	\textsc{tm}3	have.a.rest\\
\glt `and when he gets exhausted, he leaves it to have a rest.'
\z

\ea\label{ex:text18-14}
A koro no komano sinivi potee na varuruva koana, sa potee haana bono komano sinivi suuna,\\
\gll a	koro	n=o	koma-n=o	sinivi	potee	na	varuruva	koa=na	sa	potee	haa=na	bono	koma-n=o	sinivi	suuna\\
     \textsc{art}2.\textsc{sg}	carve.out	3\textsc{sg}.\textsc{poss}=\textsc{art}3.\textsc{sg}	inside-3\textsc{sg}.\textsc{poss}=\textsc{art}3.\textsc{sg}	canoe	potee	\textsc{tm}2	easy	just=3\textsc{sg}.\textsc{ipfv}	\textsc{neg}1	be.like	\textsc{neg}=3\textsc{sg}.\textsc{ipfv}	\textsc{obj}.\textsc{art}3.\textsc{sg}	inside-3\textsc{sg}.\textsc{poss}=\textsc{art}3.\textsc{sg}	canoe	indigenous\\
\glt `The carving of the inside of the \textit{potee} canoe is quite easy, it is not like the inside of the indigenous canoe,'
\z

\largerpage[2]
\ea\label{ex:text18-15}
tabae na antee koa vonaen ei bean parata o buo kubisi, o kubis beera ae o kubis bebeahu.\\
\gll tabae	na	antee	koa	vo=na=en	ei	be=an	parata	o	buo	kubisi	o	kubis-beera	ae	o	kubis-bebeahu\\
     because	\textsc{tm}	reach	only	like=3\textsc{sg}.\textsc{ipfv}=\textsc{dem}3	\textsc{dem}2	when=2\textsc{sg}.\textsc{pron}	open	\textsc{art}3.\textsc{sg}	two3.\textsc{sg}	finger	\textsc{art}3.\textsc{sg}	finger-big	and1	\textsc{art}3.\textsc{sg}	finger-long\\
\glt `because (the indigenous canoe) is only as wide as the span between your thumb and your middle finger.'\footnote{lit. `because only reaches like this: this (is) when you open the two fingers, the big finger and the long finger'}
\z\clearpage

\ea\label{ex:text18-16}
Bean boha voen paa anteu, ean repaa nao vo teo hum vaa tea pisipisi.\\
\gll be=an	boha	vo=en	paa	ante=u	e-an	re=paa	nao	vo	te=o	hum	vaa	tea	pisi{$\sim$}pisi\\
     when=2\textsc{sg}.\textsc{pron}	think	like=\textsc{dem}3	\textsc{tm}3	enough=\textsc{pfv}	\textsc{h}-2\textsc{sg}.\textsc{pron}	\textsc{consec}=\textsc{tm}3	go	\textsc{goal}	\textsc{prep}=\textsc{art}3.\textsc{sg}	part	\textsc{lnk}	\textsc{compl}1	\textsc{red}{$\sim$}carve.a.pointed.end\\
\glt `When you think that it is enough, then you go on with the part of carving the pointed ends.'
\z

\ea\label{ex:text18-17}
Bean taraha pisi vakavara a peha papana, ean repaa nao pete vo tea meha papana.\\
\gll be=an	taraha	pisi	va-kavara	a	peha	papana	e-an	re=paa	nao	pete	vo	te=a	meha	papana\\
     when=2\textsc{sg}.\textsc{pron}	plane.surface	cut.a.pointed.end	\textsc{advr}-complete	\textsc{art}2.\textsc{sg}	one2.\textsc{sg}	side	\textsc{h}-2\textsc{sg}.\textsc{pron}	\textsc{consec}=\textsc{tm}3	go	subsequently	\textsc{goal}	\textsc{prep}=\textsc{art}2.\textsc{sg}	other2.\textsc{sg}	side\\
\glt `When you have finished planing one side of the pointed end, then you go on with the other side.'
\z

\ea\label{ex:text18-18}
Be kavara, ean repaa koro vahabana o komano sinivi.\\
\gll be	kavara	e-an	re=paa	koro	vahabana	o	koma-n=o	sinivi\\
     when	finished	\textsc{h}-2\textsc{sg}.\textsc{pron}	\textsc{consec}=\textsc{tm}3	carve.out	again	\textsc{art}3.\textsc{sg}	inside-3\textsc{sg}.\textsc{poss}=\textsc{art}3.\textsc{sg}	canoe\\
\glt `After that, you carve again the inside of the canoe.'
\z

\ea\label{ex:text18-19}
Koro, koro, bean boha voen, ``antee vai,'' ean repaa nao pete teo hum vaa tea govagovara o sinivi.\\
\gll koro	koro	be=an	boha	vo=en	antee	vai	e-an	re=paa	nao	pete	te=o	hum	vaa	tea	gova{$\sim$}govara	o	sinivi\\
     carve.out	carve.out	when=2\textsc{sg}.\textsc{pron}	think	like=\textsc{dem}3	enough	now	\textsc{h}-2\textsc{sg}.\textsc{pron}	\textsc{consec}=\textsc{tm}3	go	subsequently	\textsc{prep}=\textsc{art}3.\textsc{sg}	part	\textsc{lnk}	\textsc{compl}1	\textsc{red}{$\sim$}carve.pointed.end	\textsc{art}3.\textsc{sg}	canoe\\
\glt `Carve, carve, (and) when you think, ``Enough now,'' you go on with the part of carving the bow and the stern of the canoe.'
\z

\ea\label{ex:text18-20}
Bean mei nom tea govara o sinivi, ean toro verete roho ei.\\
\gll be=an	mei	nom	tea	govara	o	sinivi	e-an	toro	verete	roho	ei\\
     when=2\textsc{sg}.\textsc{pron}	do.not.yet	\textsc{ipfv}	\textsc{compl}1	carve.pointed.end	\textsc{art}3.\textsc{sg}	canoe	\textsc{h}-2\textsc{sg}.\textsc{pron}	\textsc{tm}7	turn.over	before	\textsc{dem}2\\
\glt `Before you carve the bow and the stern of the canoe, you must turn it over.'
\z

\ea\label{ex:text18-21}
A kiu bona a beera,\\
\gll a	kiu	bona	a	beera\\
     \textsc{art}2.\textsc{sg}	work	\textsc{ana}.\textsc{sg}	\textsc{art}2.\textsc{sg}	big\\
\glt `This work is big,'
\z

\ea\label{ex:text18-22}
meo aba o sii bero vareko toro paku bari.\\
\gll me=o	aba	o	sii	bero	vareko	toro	paku	bari\\
     and4=\textsc{art}2.\textsc{pl}	person	\textsc{art}2.\textsc{pl}	\textsc{dim}	many	quite	\textsc{tm}7	do	4\textsc{pl}.\textsc{pron}\\
\glt `and more than one person must do it.'
\z

\ea\label{ex:text18-23}
Beam verete vakavara e, ean repaa taneo a kiu vaa tea govagovara.\\
\gll be=am	verete	va-kavara	e	e-an	re=paa	taneo	a	kiu	vaa	tea	gova{$\sim$}govara\\
     when=2\textsc{pl}.\textsc{pron}	turn.over	\textsc{advr}-complete	3\textsc{sg}.\textsc{pron}	\textsc{h}-2\textsc{sg}.\textsc{pron}	\textsc{consec}=\textsc{tm}3	start	\textsc{art}2.\textsc{sg}	work	\textsc{lnk}	\textsc{compl}1	\textsc{red}{$\sim$}carve.pointed.end\\
\glt `When you have turned it over, you start the work of carving the bow and the stern.'
\z

\ea\label{ex:text18-24}
Govara kaakoo teo ihu no sinivi,\\
\gll govara	kaakoo	te=o	ihu	n=o	sinivi\\
     carve.pointed.end	move	\textsc{prep}=\textsc{art}3.\textsc{sg}	nose	3\textsc{sg}.\textsc{poss}=\textsc{art}3.\textsc{sg}	canoe\\
\glt `Carve by moving towards the nose of the canoe,'
\z

\ea\label{ex:text18-25}
ean repaa vaatotobin e,\\
\gll e-an	re=paa	va-totobin	e\\
     \textsc{h}-2\textsc{sg}.\textsc{pron}	\textsc{consec}=\textsc{tm}3	\textsc{caus}-straight	3\textsc{sg}.\textsc{pron}\\
\glt `then you make it straight,'
\z

\newpage
\ea\label{ex:text18-26}
ean repaa govara pete o hum vaa teo itao.\\
\gll e-an	re=paa	govara	pete	o	hum	vaa	te=o	itao\\
     \textsc{h}-2\textsc{sg}.\textsc{pron}	\textsc{consec}=\textsc{tm}3	carve.pointed.end	subsequently	\textsc{art}3.\textsc{sg}	part	\textsc{lnk}	\textsc{prep}=\textsc{art}3.\textsc{sg}	stern\\
\glt `then you continue carving the part at the stern.'
\z

\ea\label{ex:text18-27}
Bea govara bona kavara, ean paa vaatotobin e.\\
\gll be=a	govara	bona	kavara	e-an	paa	va-totobin	e\\
     when=\textsc{art}2.\textsc{sg}	carve.pointed.end	\textsc{ana}.\textsc{sg}	finished	\textsc{h}-2\textsc{sg}.\textsc{pron}	\textsc{tm}3	\textsc{caus}-straight	3\textsc{sg}.\textsc{pron}\\
\glt `When this carving of the bow and the stern is finished, you make it straight.'
\z

\ea\label{ex:text18-28}
Be kavara, ean repaa verete vahae.\\
\gll be	kavara	e-an	re=paa	verete	vaha=e\\
     when	finished	\textsc{h}-2\textsc{sg}.\textsc{pron}	\textsc{consec}=\textsc{tm}3	turn.over	again=3\textsc{sg}.\textsc{pron}\\
\glt `When (it) is finished, then you turn it over again.'
\z

\ea\label{ex:text18-29}
O sinivi paa aeae vaevuru,\\
\gll o	sinivi	paa	aeae	vaevuru\\
     \textsc{art}3.\textsc{sg}	canoe	\textsc{tm}3	light	already\\
\glt `The canoe is already light now,'
\z

\ea\label{ex:text18-30}
me ean koa a peha na antee nom tea verete vahae.\\
\gll me	e-an	koa	a	peha	na	antee	nom	tea	verete	vaha=e\\
     and4	\textsc{h}-2\textsc{sg}.\textsc{pron}	only	\textsc{art}2.\textsc{sg}	one2.\textsc{sg}	\textsc{tm}2	can	\textsc{ipfv}	\textsc{compl}1	turn.over	again=3\textsc{sg}.\textsc{pron}\\
\glt `and you can turn it over again just by yourself.'\footnote{lit. `and you only, (being) one, can turn it over again.'}
\z

\ea\label{ex:text18-31}
Bean verete vakavara, ean repaa koro vahabana e.\\
\gll be=an	verete	va-kavara	e-an	re=paa	koro	vahabana	e\\
     when=2\textsc{sg}.\textsc{pron}	turn.over	\textsc{advr}-complete	\textsc{h}-2\textsc{sg}.\textsc{pron}	\textsc{consec}=\textsc{tm}3	carve.out	again	3\textsc{sg}.\textsc{pron}\\
\glt `When you have turned it over, then you do some more carving on the inside.'
\z

\ea\label{ex:text18-32}
Koro, koro, bean boha voen paa anteu, ean paa antee,\\
\gll koro	koro	be=an	boha	vo=en	paa	ante=u	e-an	paa	antee\\
     carve.out	carve.out	when=2\textsc{sg}.\textsc{pron}	think	like=\textsc{dem}3	\textsc{tm}3	enough=\textsc{pfv}	\textsc{h}-2\textsc{sg}.\textsc{pron}	\textsc{tm}3	stop\\
\glt `Carve (and) carve, (and) when you think that it is enough, you stop,'
\z

\ea\label{ex:text18-33}
ei ihuana vai o hum vaa tea rahi voosu.\\
\gll ei	ihuana	vai	o	hum	vaa	tea	rahi	voosu\\
     \textsc{dem}2	wait	now	\textsc{art}3.\textsc{sg}	part	\textsc{lnk}	\textsc{compl}1	pull	go.home\\
\glt `(and) now (you) wait for the part of pulling (the canoe) home.'
\z

\ea\label{ex:text18-34}
Teo bon rahi no potee, ean repaa sue ki to sii aba ore ausu vuan tea rahi voosu e.\\
\gll te=o	bon	rahi	n=o	potee	e-an	re=paa	sue	ki	to	sii	aba	o=re	ausu	vu=an	tea	rahi	voosu	e\\
     \textsc{prep}=\textsc{art}3.\textsc{sg}	day	pull	3\textsc{sg}.\textsc{poss}=\textsc{art}3.\textsc{sg}	potee.canoe	\textsc{h}-2\textsc{sg}.\textsc{pron}	\textsc{consec}=\textsc{tm}3	say	\textsc{dat}	\textsc{nspec}2.\textsc{pl}	\textsc{dim}	person	3.\textsc{pron}=\textsc{consec}	assist	2\textsc{sg}.\textsc{objm}=2\textsc{sg}.\textsc{pron}	\textsc{compl}1	pull	go.home	3\textsc{sg}.\textsc{pron}\\
\glt `As for the day of pulling the \textit{potee} canoe, you tell some people to assist you to pull it home.'
\z

\ea\label{ex:text18-35}
A kiu vai ahiki ta von,\\
\gll a	kiu	vai	ahiki	ta	von\\
     \textsc{art}2.\textsc{sg}	work	\textsc{dem}6	not.exist	\textsc{nspec}2.\textsc{sg}	payment\\
\glt `This work does not cost anything,'\footnote{lit. `This work, there isn't any payment.'}
\z

\ea\label{ex:text18-36}
evehee ean toro paku a sii tabaan to kahi ma an koa vue o manae to paa nao nau tea rahi sinivi.\\
\gll evehee	e-an	toro	paku	a	sii	tabaan	to	kahi	ma	an	koa	vue	o	manae	to	paa	nao	na=u	tea	rahi	sinivi\\
     but	\textsc{h}-2\textsc{sg}.\textsc{pron}	\textsc{tm}7	make	\textsc{art}2.\textsc{sg}	\textsc{dim}	food	\textsc{rel}	\textsc{tm}1	come	eat	just	particulary	\textsc{art}2.\textsc{pl}	young.men	\textsc{rel}	\textsc{tm}3	go	thither=\textsc{pfv}	\textsc{compl}1	pull	canoe\\
\glt `but you should make some food that the young men will eat, who have gone there to pull the canoe.'
\z

\ea\label{ex:text18-37}
Beori rahi voosu maa bari, ean repaa vaan riori, a menaga, nahu iana bara ta maa meha tabaan bata mee.\\
\gll be=ori	rahi	voosu	maa	bari	e-an	re=paa	va-an	ri=ori	a	menaga	nahu	iana	bara	ta	maa	meha	tabaan	bata	mee\\
     when=3\textsc{pl}.\textsc{pron}	pull	go.home	hither	4\textsc{pl}.\textsc{pron}	\textsc{h}-2\textsc{sg}.\textsc{pron}	\textsc{consec}=\textsc{tm}3	\textsc{caus}-eat	3\textsc{pl}.\textsc{objm}=3\textsc{pl}.\textsc{pron}	\textsc{art}2.\textsc{sg}	manioc.dumplings	pot	fish	and2	\textsc{nspec}2.\textsc{sg}	\textsc{plm}	other2.\textsc{sg}	food	\textsc{cont}	also\\
\glt `When they have pulled it home, you will give them to eat manioc dumplings, a pot of fish, and also some other food along with it.'
\z

\ea\label{ex:text18-38}
Bea maa taba vai kavara, ean paa paku papahiana komana nom amaa sii vagegeana kiu.\\
\gll be=a	maa	taba	vai	kavara	e-an	paa	paku	papahiana	komana	nom	a=maa	sii	vagegeana	kiu\\
     when=\textsc{art}2.\textsc{sg}	\textsc{plm}	thing	\textsc{dem}6	finished	\textsc{h}-2\textsc{sg}.\textsc{pron}	\textsc{tm}3	do	alone	self	\textsc{ipfv}	\textsc{art}2.\textsc{sg}=\textsc{plm}	\textsc{dim}	final	work\\
\glt `When these things are finished, you will do the final little pieces of work just by yourself.'
\z

\ea\label{ex:text18-39}
Beo naono tei me nana bona maa panapanaa, ean repaa gono maa ta maa vua asita tea asita o sinivi tean.\\
\gll be=o	naono	tei	me	nana	bona	maa	panapanaa	e-an	re=paa	gono	maa	ta	maa	asita	tea	asita	o	sinivi	te=an\\
     when=\textsc{art}3.\textsc{sg}	wood	be.there	with	3\textsc{sg}.\textsc{ipfv}	\textsc{obj}.\textsc{art}2.\textsc{sg}	\textsc{plm}	knothole	\textsc{h}-2\textsc{sg}.\textsc{pron}	\textsc{consec}=\textsc{tm}3	get	hither	\textsc{nspec}2.\textsc{sg}	\textsc{plm}	putty.nut	\textsc{compl}1	seal.with.putty.nut.glue	\textsc{art}2.\textsc{sg}	canoe	\textsc{prep}=2\textsc{sg}.\textsc{pron}\\
\glt `When the wood has knotholes, you get some putty nuts to seal your canoe with putty nut glue.'
\z

\ea\label{ex:text18-40}
A meha maa naono na hiki ni nana bona panapanaa.\\
\gll a	meha	maa	naono	na	hiki	ni	nana	bona	panapanaa\\
     \textsc{art}2.\textsc{sg}	other2.\textsc{sg}	\textsc{plm}	tree	\textsc{tm}2	not.have	\textsc{appl}	3\textsc{sg}.\textsc{ipfv}	\textsc{obj}.\textsc{art}3.\textsc{pl}	knothole\\
\glt `Some trees do not have any knotholes.'
\z

\newpage
\ea\label{ex:text18-41}
Eikou ean repaa pahin paku bata koa amaa tapara.\\
\gll eikou	e-an	re=paa	pahin	paku	bata	koa	a=maa	tapara\\
     and.so	\textsc{h}-2\textsc{sg}.\textsc{pron}	\textsc{consec}=\textsc{tm}3	straightaway	make	\textsc{cont}	just	\textsc{art}2.\textsc{sg}=\textsc{plm}	tapara.pole\\
\glt `And so you straightaway make some \textit{tapara} poles (that prevent the sides from moving inside).'
\z

\ea\label{ex:text18-42}
O naono vai to tei me nana bona panapanaa toro vaatei roho beo asita raka roho nana.\\
\gll o	naono	vai	to	tei	me	nana	bona	panapanaa	toro	vaa-tei	roho	be=o	asita	raka	roho	nana\\
     \textsc{art}3.\textsc{sg}	wood	\textsc{dem}6	\textsc{rel}	be.there	with	3\textsc{sg}.\textsc{ipfv}	\textsc{obj}.\textsc{art}3.\textsc{pl}	knothole	\textsc{tm}7	\textsc{caus}-be.there	before	when=\textsc{art}3.\textsc{sg}	putty.nut.glue	dry	before	3\textsc{sg}.\textsc{ipfv}.\\
\glt `The wood that has knotholes must be left (untouched) until the glue has dried up.'
\z

\ea\label{ex:text18-43}
Bean paku vakavara amaa tapara, ean repaa nao vo teo hum vaa tea pam no sinivi.\\
\gll be=an	paku	va-kavara	a=maa	tapara	e-an	re=paa	nao	vo	te=o	hum	vaa	te=a	pam	n=o	sinivi\\
     when=2\textsc{sg}.\textsc{pron}	make	\textsc{advr}-complete	\textsc{art}2.\textsc{sg}=\textsc{plm}	tapara.pole	\textsc{h}-2\textsc{sg}.\textsc{pron}	\textsc{consec}=\textsc{tm}3	go	\textsc{goal}	\textsc{prep}=\textsc{art}3.\textsc{sg}	part	\textsc{lnk}	\textsc{prep}=\textsc{art}2.\textsc{sg}	tie.up	3\textsc{sg}.\textsc{poss}=\textsc{art}3.\textsc{sg}	canoe\\
\glt `When you have finished fixing the \textit{tapara} poles, you move on to the part of tying up the canoe.'
\z

\ea\label{ex:text18-44}
Teo potee, o buaku koa o vanai, o peho matana, o meho he itao.\\
\gll te=o	potee	o	buaku	koa	o   vanai	o	peho	matana	o	meho	he	itao\\
     \textsc{prep}=\textsc{art}3.\textsc{sg}	potee.canoe	\textsc{art}3.\textsc{sg}	two	only	\textsc{art}3.\textsc{sg}	outrigger.boom	\textsc{art}3.\textsc{sg}	one3.\textsc{sg}	bow	\textsc{art}3.\textsc{sg}	other3.\textsc{sg}	\textsc{advers}	stern\\
\glt `As for the \textit{potee} canoe, there are only two outrigger booms, one near the bow, and the other near the stern.'
\z

\newpage
\ea\label{ex:text18-45}
Bean pam vakavara nio buo vanai, ean repaa tapa ae pam vagegeana koa nio aga ore paa kavara.\\
\gll be=an	pam	va-kavara	ni=o	buo	vanai	e-an	re=paa	tapa	ae	pam	vagegeana	koa	ni=o	aga	o=re	paa	kavara\\
     when=2\textsc{sg}.\textsc{pron}	tie.up	\textsc{advr}-complete	\textsc{appl}=\textsc{art}3.\textsc{sg}	two3.\textsc{sg}	outrigger.boom	\textsc{sg}.\textsc{pron}	\textsc{consec}=\textsc{tm}3	fix.with.pegs	and4	tie.up	finally	just	\textsc{appl}=\textsc{art}3.\textsc{sg}	float	3.\textsc{pron}=\textsc{consec}	\textsc{tm}3	finished\\
\glt `When you have tied the two outrigger booms (to the hull of the canoe), you finally fix the float with pegs and lashings (onto the outrigger booms), and then it is finished.'
\z

\ea\label{ex:text18-46}
Ei ean antee vai e tea hae tea vagana ge havee.\\
\gll ei	e-an	antee	vai	e	tea	hae	tea	vagana	ge	havee\\
     \textsc{dem}2	\textsc{h}-2\textsc{sg}.\textsc{pron}	can	now	3\textsc{sg}.\textsc{pron}	\textsc{compl}1	board	\textsc{compl}1	fish	or	what\\
\glt `Now you can board it to go fishing or whatever.'
\z

\ea\label{ex:text18-47}
O kaku sinivi vai o sinivi vaa Solomon.\\
\gll o	kaku	sinivi	vai	o	sinivi	vaa	solomon\\
     \textsc{art}3.\textsc{sg}	type	canoe	\textsc{dem}6	\textsc{art}3.\textsc{sg}	canoe	\textsc{lnk}	Solomon\\
\glt `This type of canoe is a canoe from the Solomon Islands.'
\z

\ea\label{ex:text18-48}
Evehee a ta Solomon saka vaatei haari bono aga.\\
\gll evehee	a	ta	solomon	saka	vaa-tei	haa=ri	bono	aga\\
     but	\textsc{art}2.\textsc{sg}	people	Solomon	\textsc{neg}1	\textsc{caus}-be.there	\textsc{neg}2=3\textsc{pl}.\textsc{ipfv}	\textsc{obj}.\textsc{art}3.\textsc{sg}	float\\
\glt `But the Solomon people do not attach a float.'
\z

\ea\label{ex:text18-49}
Eve he nomaa voen meara paa vaatei o aga,\\
\gll eve	he	nomaa	vo=en	me-ara	paa	vaa-tei	o	aga\\
     \textsc{h}.3\textsc{sg}.\textsc{pron}	\textsc{advers}	come	like=\textsc{dem}3	and4-1\textsc{pl}.\textsc{in}.\textsc{pron}	\textsc{tm}3	\textsc{caus}-be.there	\textsc{art}3.\textsc{sg}	float\\
\glt `When it came (to us) like this, we attached the float,'
\z

\newpage
\ea\label{ex:text18-50}
ei kou a tahii teara na pohapoha sana nana.\\
\gll ei~kou	a	tahii	te=ara	na	poha{$\sim$}poha	sana	nana\\
     because	\textsc{art}2.\textsc{sg}	sea	\textsc{prep}=1\textsc{pl}.\textsc{in}.\textsc{pron}	\textsc{tm}2	\textsc{red}{$\sim$}wave	very	3\textsc{sg}.\textsc{ipfv}\\
\glt `because our sea is very rough.'
\z

\ea\label{ex:text18-51}
O potee na antee nana tea tatana bona too hata.\\
\gll o	potee	na	antee	nana	tea	tatana	bona	too	hata\\
     \textsc{art}3.\textsc{sg}	potee.canoe	\textsc{tm}2	can	3\textsc{sg}.\textsc{ipfv}	\textsc{compl}1	resist	\textsc{obj}.\textsc{art}2.\textsc{sg}	weather	bad\\
\glt `The \textit{potee} canoe can resist bad weather.'
\z

\ea\label{ex:text18-52}
Sa manin goo haana.\\
\gll sa	manin	goo	haa=na\\
     \textsc{neg}1	ever	sink	\textsc{neg}2=3\textsc{sg}.\textsc{ipfv}\\
\glt `It never sinks,'
\z

\ea\label{ex:text18-53}
o sinivi suuna he na antee nana tea goo.\\
\gll o	sinivi	suuna	he	na	antee	nana	tea	goo\\
     \textsc{art}3.\textsc{sg}	canoe	indigenous	\textsc{advers}	\textsc{tm}2	can	3\textsc{sg}.\textsc{ipfv}	\textsc{compl}1	sink\\
\glt `whereas the indigenous canoe can sink.'
\z
